%% ============================================================
%%  SECCIÓN II — Resumen
%% ============================================================

\section{Resumen}

Se estudió experimentalmente el comportamiento elástico de dos materiales sometidos a esfuerzos de tensión: un resorte metálico helicoidal y una tira de jebe (liga elastomérica). Para cada material se aplicaron seis cargas acumulativas en el rango de \qtyrange{500}{1750}{g}, registrando la deformación longitudinal y las dimensiones de la sección transversal tanto en la fase de carga como en la de descarga. A partir de estos datos se calcularon el esfuerzo real $\sigma = F/S$ (considerando la reducción dinámica del área transversal) y la deformación unitaria $\varepsilon = \Delta\ell/\ell_0$, construyendo los diagramas $F$\,vs.\,$\Delta\ell$ y $\sigma$\,vs.\,$\varepsilon$ para caracterizar su respuesta mecánica.

El resorte metálico exhibió una relación lineal entre fuerza y deformación dentro del rango ensayado ($R^2 = 0{,}9966$), verificando la validez de la Ley de Hooke. El ajuste por mínimos cuadrados entregó una constante recuperadora $k = \qty{80,89 \pm 2,38}{N/m}$ y un módulo de Young efectivo recalculado de $Y_r = \qty{90,23 \pm 2,41}{kPa}$. El trabajo elástico conservativo realizado hasta la tercera carga resultó $W = \qty{0,715}{J}$. 

La liga, en cambio, presentó un comportamiento marcadamente no lineal con endurecimiento progresivo, efecto amplificado al utilizar el esfuerzo real en el análisis. Se obtuvo un módulo secante promedio de $Y_l = \qty{538 \pm 25}{kPa}$. Asimismo, el material exhibió una notable histéresis elástica; la integración numérica del lazo en el diagrama de esfuerzo real permitió cuantificar una densidad de energía disipada por unidad de volumen de $A = \qty{13,817}{kJ/m^3}$. La comparación entre ambos ilustra la diferencia entre un sólido elástico lineal y un elastómero viscoelástico, cuya respuesta disipa energía térmica en cada ciclo de deformación.